\section{Related Work}
\label{sec:related}
\subsection{Dynamic Scene Reconstruction}
Recent works have leveraged NeRF representations to jointly solve for canonical space and deformation fields in dynamic scenes using RGB supervision~\citep{guo2023forward, li2021neural, park2021nerfies, park2021hypernerf, pumarola2020dnerf, tretschk2021nonrigid, xian2021spacetime,liu2023robust,wu2024denver,chen2024narcan,ma2024humannerf}. Further advancements in dynamic neural rendering include object segmentation~\citep{song2023nerfplayer}, incorporation of depth information~\citep{attal2021torf}, utilization of 2D CNNs for scene priors~\citep{lin2022efficient, peng2023representing}, and multi-view video compression~\citep{li2022neural3D}. However, these NeRF-based methods are computationally intensive, limiting their practical applications. To address this, 3D Gaussian Splatting~\citep{kerbl20233d} has emerged as a promising alternative, offering real-time rendering capabilities while maintaining high visual quality. Building upon this efficient representation, recent research has adapted 3D Gaussians for dynamic scenes~\citep{yang2023deformable, wu20234d, huang2024scgs, liang2023gaufre,wang2024shape,mihajlovic2024splatfields,stearns2024dynamic}. Nevertheless, these approaches do not explicitly account for changes in surface normal during the dynamic process. Our work extends this line of research by combining specular object rendering based on normal estimation with a deformation field, enabling each 3D Gaussian to model dynamic specular scenes effectively.

\subsection{Reflective Object Rendering}
While significant progress has been made in rendering reflective objects, challenges from complex light interactions persist. Recent years have seen numerous studies addressing these issues, primarily by decomposing appearance into lighting and material properties ~\citep{bi2020neural, boss2021neuralpil,li2022neural, srinivasan2020nerv, Zhang_2021, Munkberg_2022_CVPR,zhang2021physg,verbin2024eclipse,zhao2024illuminerf}. Building on this foundation, some research has focused on improving the capture and reproduction of specular reflections ~\citep{verbin2022ref, ma2023specnerf, verbin2024nerf, ye20243d}. In contrast, others have leveraged signed distance functions (SDFs) to enhance normal estimation ~\citep{ge2023refneus, liang2023envidr, liang2023spidr, liu2023nero, zhang2023neilf}.
The emergence of 3D Gaussian Splatting (3DGS) has sparked a new wave of techniques ~\citep{jiang2023gaussianshader,liang2023gs,yang2024specgaussian,zhu2024gs,gao2023relightable,shi2023gir} that integrate Gaussian splatting with physically-based rendering. Nevertheless, accurately modeling dynamic environments and time-varying specular reflections remains a significant challenge. To address this limitation, our work introduces a novel approach incorporating a deformable environment map and additional explicit Gaussian attributes specifically designed to capture specular color changes over time.
